% ============================================================
%  Embodied Intelligence Survey
%  Compiled with XeLaTeX
%  Font: Times New Roman, 12pt (小四), 1.5× line spacing
% ============================================================
\documentclass[12pt, a4paper]{article}

% ---------- encoding & font ----------
\usepackage{fontspec}
\setmainfont{Times New Roman}

% ---------- layout ----------
\usepackage[
  top=2.54cm, bottom=2.54cm,
  left=3.17cm, right=3.17cm
]{geometry}
\usepackage{setspace}
\onehalfspacing                  % 1.5× line spacing

% ---------- math ----------
\usepackage{amsmath, amssymb}

% ---------- graphics & tables ----------
\usepackage{graphicx}
\usepackage{booktabs}
\usepackage{caption}
\captionsetup{font=small, labelfont=bf}

% ---------- hyperlinks ----------
\usepackage[
  colorlinks=true,
  linkcolor=black,
  citecolor=blue,
  urlcolor=blue
]{hyperref}

% ---------- bibliography ----------
\usepackage[numbers, sort&compress]{natbib}
\bibliographystyle{unsrtnat}     % numbered, in citation order

% ---------- miscellaneous ----------
\usepackage{enumitem}
\usepackage{xcolor}

% ============================================================
\begin{document}

% ---- Title block ----
\title{\textbf{Embodied Intelligence: A Survey of Foundations,\\
               Learning Paradigms, and Open Challenges}}

\author{
  Your Name\textsuperscript{1}
  \\[4pt]
  \small\textsuperscript{1}Department of XXX, University of XXX, City, Country\\
  \small\texttt{your.email@university.edu}
}

\date{\today}
\maketitle

% ============================================================
\begin{abstract}
% TODO (write last): 250–300 words covering scope, taxonomy,
%                    key findings, and open challenges.
Embodied intelligence refers to the study of intelligent systems
that perceive, reason, and act through a physical body embedded
in a real or simulated environment.
Unlike classical disembodied AI, embodied agents must
continuously interact with their surroundings, making
sensorimotor coordination, grounded language understanding,
and adaptive planning central research concerns.
This survey provides a structured overview of the field,
covering: (i)~the foundational definitions and historical
development of embodied cognition; (ii)~core learning
paradigms including reinforcement learning, imitation learning,
and vision-language-action (VLA) models; (iii)~representative
tasks and benchmarks; (iv)~evaluation protocols and datasets;
and (v)~outstanding open challenges.
We identify promising future directions such as world-model-based
planning, scalable multi-embodiment transfer, and
human–robot collaborative learning.

% TODO: replace placeholder text with your own summary after
%       the main sections are drafted.

\noindent\textbf{Keywords:} embodied intelligence, embodied AI,
reinforcement learning, imitation learning,
vision-language-action models, robot learning, world models
\end{abstract}

% ============================================================
\section{Introduction}
\label{sec:intro}

% --- Background ---
Intelligent agents that operate through physical bodies occupy
a central position in both biological intelligence and
artificial intelligence research.
The hypothesis that ``intelligence requires a body'' was
articulated by Moravec~\cite{moravec1988mind} and later
formalised under the banner of \emph{embodied cognition}
by Brooks~\cite{brooks1991intelligence}, who argued that
cognition emerges from the dynamic coupling between an
agent and its environment rather than from disembodied
symbol manipulation.

% --- Why now ---
Recent advances in deep learning, large language models (LLMs),
and robot hardware have created an unprecedented opportunity to
revisit these ideas at scale.
Foundation models pre-trained on internet-scale data can be
adapted to robotic control~\cite{brohan2022rt1, driess2023palme},
and simulated environments now provide photorealistic training
grounds with efficient physics simulation~\cite{kolve2017ai2thor,
savva2019habitat}.

% --- Scope of this survey ---
This survey is organised as follows.
Section~\ref{sec:background} introduces the conceptual
foundations and key terminology.
Section~\ref{sec:methods} reviews the major learning paradigms.
Section~\ref{sec:tasks} surveys benchmark tasks and datasets.
Section~\ref{sec:challenges} discusses remaining open challenges.
Section~\ref{sec:outlook} provides a summary and outlook.

% ============================================================
\section{Background and Foundations}
\label{sec:background}

\subsection{Definition and Scope}
% TODO: define embodied intelligence; contrast with classical AI
% Key citations: Brooks 1991, Pfeifer & Scheier 1999,
%                Anderson 2003 (embodied cognition review)

\subsection{Historical Development}
% TODO: cybernetics -> behaviour-based robotics ->
%       deep RL era -> foundation model era
% Key citations: Wiener 1948, Brooks 1986 subsumption,
%                Mnih 2015 DQN, Brohan 2022 RT-1

\subsection{Key Concepts}
% TODO: grounding, embodiment, closed-loop control,
%       perception-action loop, world model

% ============================================================
\section{Learning Paradigms}
\label{sec:methods}

\subsection{Reinforcement Learning for Embodied Agents}
% TODO: model-free (PPO, SAC), model-based (Dreamer, MBPO),
%       sim-to-real transfer
% Key citations: Mnih 2015, Schulman 2017 PPO,
%                Haarnoja 2018 SAC, Hafner 2019 Dreamer

\subsection{Imitation Learning and Learning from Demonstration}
% TODO: BC, GAIL, DAgger, diffusion policy
% Key citations: Ho 2016 GAIL, Ross 2011 DAgger,
%                Chi 2023 Diffusion Policy, Jang 2022 BC-Z

\subsection{Vision-Language-Action Models}
% TODO: RT-1, RT-2, PaLM-E, SayCan, OpenVLA
% Key citations: Brohan 2022 RT-1, Brohan 2023 RT-2,
%                Driess 2023 PaLM-E, Ahn 2022 SayCan

\subsection{World Models and Planning}
% TODO: model-based planning, Dreamer, UniSim, GAIA
% Key citations: Hafner 2020 Dreamer v2, Yang 2023 UniSim

\subsection{Multi-modal Perception}
% TODO: visual grounding, depth, tactile, proprioception
% Key citations: Radford 2021 CLIP, Gao 2023 ...

% ============================================================
\section{Benchmark Tasks and Datasets}
\label{sec:tasks}

\subsection{Navigation and Exploration}
% TODO: PointNav, ObjectNav, VLN
% Key citations: Anderson 2018 R2R, Savva 2019 Habitat,
%                Kolve 2017 AI2-THOR

\subsection{Manipulation}
% TODO: pick-and-place, dexterous grasping, deformable objects
% Key citations: Zhu 2020 robosuite, Yu 2020 Meta-World,
%                Mandlekar 2021 robomimic

\subsection{Instruction Following and Rearrangement}
% TODO: ALFRED, TEACh, BEHAVIOR-1K
% Key citations: Shridhar 2020 ALFRED, Francis 2023 BEHAVIOR-1K

\subsection{Open-World and Long-Horizon Tasks}
% TODO: MineDojo, Voyager, GROOT
% Key citations: Fan 2022 MineDojo, Wang 2023 Voyager

% ============================================================
\section{Open Challenges}
\label{sec:challenges}

\begin{enumerate}[leftmargin=*, label=\textbf{C\arabic*.}]

  \item \textbf{Sim-to-Real Transfer.}
  % TODO: domain randomisation, adaptive policies,
  %       photorealistic simulation

  \item \textbf{Sample Efficiency.}
  % TODO: data augmentation, offline RL, pre-training

  \item \textbf{Generalisation Across Embodiments.}
  % TODO: morphology-agnostic policies, multi-robot transfer

  \item \textbf{Safety and Alignment.}
  % TODO: safe RL, constraint satisfaction, human oversight

  \item \textbf{Long-Horizon Reasoning.}
  % TODO: hierarchical RL, task-and-motion planning,
  %       LLM-based subgoal generation

  \item \textbf{Evaluation Standardisation.}
  % TODO: fragmented benchmarks, reproducibility

\end{enumerate}

% ============================================================
\section{Summary and Outlook}
\label{sec:outlook}

% TODO: 3–4 paragraphs
% Paragraph 1: recap what the survey covered
% Paragraph 2: most promising near-term directions
%              (VLA scaling, world models, dexterous manipulation)
% Paragraph 3: longer-term vision (AGI robots, collaborative AI)

Embodied intelligence stands at an exciting inflection point.
The convergence of large foundation models, scalable simulators,
and increasingly capable robot hardware suggests that many of the
challenges surveyed in Section~\ref{sec:challenges} may see
substantial progress in the coming years.
We anticipate that \emph{world-model-based planning},
\emph{cross-embodiment transfer}, and
\emph{human-in-the-loop learning} will be the most productive
research directions in the near term.

% ============================================================
% Bibliography
\bibliography{references}

\end{document}
